% Part: many-valued-logic
% Chapter: infinite-valued-logics
% Section: introduction

\documentclass[../../../include/open-logic-section]{subfiles}

\begin{document}

\olfileid{mvl}{inf}{int}

\olsection{పరిచయం}

ఒక మాత్రికలో సత్యమూల్యాల సంఖ్య పరిమితమైనదే
కానక్కర్లేదు. అనంత సత్యమూల్యాల సమితికి
సహజమైన ఎంపిక $0$ నుంచి~$1$ వరకు
ఉన్న కరణీయ సంఖ్యల సమితి
$V_\infty = [0,1] \cap \Rat$; అంటే
\begin{align*}
    V_\infty & = \Setabs{\frac{n}{m}}{n,m \in \Nat \text{ మరియు } 0<m \text{ మరియు } n\le m}.
\intertext{ఈ అనంత సత్యమూల్యాల సమితిని పరిగణించినప్పుడు,
రెండు లేదా అంతకన్నా ఎక్కువ విలువల కోసం కింది ఉపసమితులనూ
పరిగణించడం తరచూ ఉపయోగకరం:}
V_m & = \Setabs{\frac{n}{m-1}}{n \in \Nat \text{ మరియు } n<m}
\intertext{ఉదాహరణకు, $V_5$లో సమాన అంతరాలున్న $5$ సత్యమూల్యాలు ఉంటాయి:}
V_5 & = \{0, \frac{1}{4}, \frac{1}{2}, \frac{3}{4}, 1\}.
\end{align*}
\sourcecorrection{OLTEMVLINF-001}{ఆంగ్ల మూలంలోని
మొదటి సమితి నిర్వచనం హారం సున్నా కావడాన్ని
నిషేధించలేదు; సున్నాతో భాగించడం నిర్వచితం
కాదు. కరణీయ అంతర సమితితో సరిపడేలా
హారం సున్నాకన్నా ఎక్కువ అనే షరతు చేర్చాం.}
\sourcecorrection{OLTEMVLINF-002}{రెండవ మూల
సమితిలో లవం విలువల సంఖ్యకు సమానమయ్యే
సంభవాన్నీ అనుమతించింది. అప్పుడు అదనపు,
ఒకటికన్నా పెద్ద విలువ వచ్చి ఐదు-విలువల
ఉదాహరణకు విరుద్ధమవుతుంది. లవం విలువల
సంఖ్యకన్నా చిన్నదిగా మార్చి, కనీసం రెండు
విలువల సందర్భమని గద్యంలో స్పష్టం చేశాం.}
ఇలాంటి సత్యమూల్యాల సమితులపై ఆధారపడే
తర్కాల్లో సాధారణంగా $1$ మాత్రమే
నిర్దేశిత విలువ; అంటే $V^+ = \{1\}$.
మరో మాటలో, $1$ సంపూర్ణ సత్యానికి,
$0$ సంపూర్ణ అసత్యానికి సూచికలు;
అయితే \tetoken{సూత్రాలు}{!!{formula}s}
~$V$లోని మధ్యంతర విలువలనూ పొందవచ్చు.

$0$ నుంచి~$1$ వరకు ఉన్న అన్ని
\emph{వాస్తవ} సంఖ్యల సమితి
$V_{[0,1]} = [0,1]$ను, లేదా $[0,1]$లోని
ఇతర అనంత ఉపసమితులనూ పరిగణించవచ్చు.
ఈ సత్యమూల్యాల సమితిని వాడే తర్కాలను
తరచూ \emph{ఫజీ} తర్కాలు అంటారు.

\end{document}
